\begin{textsample}{POS Dim 4 – sc – Score 11.00 – sc/sc\_em\_54.txt}  \label{ex:f4_pos_032}
2.3 TRILHA 1 - \textbf{MULHERES} NO \textbf{TERRITÓRIO} CATARINENSE Tema : Representações sociais, \textbf{mulheres}, direitos, relações de poder e de gênero.
Área do Conhecimento : Ciências Humanas e Sociais Aplicadas.
Carga horária : 160 h ou 240 h, a depender da matriz curricular em funcionamento na Unidade Escolar.
Aulas semanais : 10 ou 15 aulas, conforme matriz da escola.
Perfil do professor : - Habilitação em Filosofia : 2 aulas ou 3 aulas, a depender da matriz curricular em funcionamento na Unidade Escolar. - Habilitação em Sociologia : 3 aulas ou 4 aulas, a depender da matriz curricular em funcionamento na Unidade Escolar. - Habilitação em História : 3 aulas ou 4 aulas, a depender da matriz curricular em funcionamento na Unidade Escolar. - Habilitação em Geografia : 2 aulas ou 3 aulas, a depender da matriz curricular em funcionamento na Unidade Escolar. 2.3.1 Unidades curriculares - Unidade curricular I - Representações sociais das \textbf{mulheres} – 40 h ou 60 h, a depender da matriz curricular implementada na Unidade Escolar. - Unidade curricular II - Mulheres nas violências de gênero – 40 h ou 60 h, a depender da matriz curricular implementada na Unidade Escolar. - Unidade curricular III - Direitos e cidadania ativa das \textbf{mulheres} - 40 h ou 60 h, a depender da matriz curricular implementada na Unidade Escolar. - Unidade curricular IV - Protagonismos das \textbf{mulheres} – 40 h ou 60 h, a depender da matriz curricular implementada na Unidade Escolar. 2.3.2 Introdução A trilha “ Mulheres no Território Catarinense ” informa haver inúmeras \textbf{mulheres} no mundo, como no Brasil e em Santa Catarina, que alcançaram seus objetivos e realizaram, neste século XXI, ações nunca antes imaginadas por gerações que viveram nos séculos passados.
Elas ocupam campos socialmente representados como exclusivos de homens ; tomam lugares de representatividade no público e no privado, ganham prêmios, lutam, resistem, vivem e \textbf{sobrevivem} na atualidade.
Neste século XXI, há um cenário inédito, com a emergência de tantas outras \textbf{mulheres} - negras, \textbf{indígenas}, caboclas, \textbf{velhas} e jovens -, conforme pesquisas publicadas na obra “ Nova História das Mulheres no Brasil ”, organizada por Carla Bassanezi Pinsky e Joana Maria Pedro ( PINSKY ; PEDRO, 2018 ).
Alicerçado nesses argumentos, o tema é instigante para os(as) estudantes do Novo Ensino Médio.
De outra parte, esse tema é duplamente justificado, tanto pelas realizações das \textbf{mulheres}, quanto porque se registram, no século XXI, desigualdades sociais graves entre homens e \textbf{mulheres} e entre \textbf{mulheres}, o que, de certa forma, explica os movimentos sociais, em específico o das \textbf{mulheres} e o das feministas, a prosseguirem lutas com pautas ainda ‘ antigas ’ e outras, mais ‘ atuais ’ ( PINTO, 2003 ).
É a violência contra as \textbf{mulheres} a faceta cruel das relações assimétricas entre esses sujeitos.
Nos campos da História e do Direito, sabe -se que elas são as vítimas de hostilidades, que vão desde a agressão ao assassinato, passando pela doméstica e a familiar, que se apresentam sob a forma de abusos sexuais, físicos, psicológicos, patrimoniais e morais, ocorridas em lugares de vivência, no trabalho ou em relações afetivas ( LAGE ; NADER, 2018 ).
Na persistência histórica da violência contra as \textbf{mulheres}, Carneiro ( 2011 ) alerta acerca de dados consistentes do desrespeito aos direitos humanos, como os apontados nos graves indicadores sociais da desigualdade de gênero e \textbf{racial}.
O mesmo faz Soihet ( 1989 ), que investigou as múltiplas formas de violência ligadas às condições sociais das \textbf{mulheres} da classe popular.
Grossi atesta da complexidade sobre as violências domésticas, que têm “ como pressuposto o fato de que não se pode isolar o polo \textbf{mulher} para entender o uso da violência em uma relação afetivo-conjugal ” ( 2011, p. 122 ).
Partindo destes estudos para aprofundar conhecimentos, é necessário visibilizar as representações sociais dessas múltiplas \textbf{mulheres}.
Para isso, utiliza -se a categoria “ mulheres ” no plural, por se entender que há sujeitos históricos diferentes em suas lutas e reivindicações ( DUBY, 1995 ).
As representações sociais, conforme explica Chartier ( 1988 ), são determinadas por interesses de grupos que as forjam.
Faz -se necessário analisar as posições de quem utiliza estas representações sociais, para exercer suas relações de poder no âmbito cultural e social.
Portanto, os sujeitos produtores e receptores de cultura estão situados entre as representações ( modos de ver ) e práticas ( modos de fazer ), cabendo -lhes analisar a posição que ocupam na rede social, como, para citar um exemplo, o da persistência histórica das violências afetivo-conjugais em nossa sociedade.
Esta trilha oferece um espaço educativo para que os(as) estudantes possam investigar os motivos pelos quais as práticas ou relações de poder engendram as representações sociais, consideradas legítimas, de comportamentos femininos e masculinos, tendo em vista a constituição social, cultural e histórica.
Usa -se a perspectiva de gênero, como ferramenta de diálogo entre campos, na área de Ciências Humanas e Sociais Aplicadas.
Esta perspectiva funciona como integradora de saberes da trilha, pois sua apropriação é usada para desnaturalizar as identidades sexuais.
Somada ao debate provocado por Pedro ( 2005 ), a categoria gênero no Brasil é útil, assim como o são as contribuições de Scott ( 1991 ), que lançou luz sobre gênero, situando -o no âmbito das relações de poder.
Nestes termos, Foucault ( 1984, 1995 ) compreende o poder pensado como positivo e como prática social.
Em outros termos, de acordo com o pensador, o poder é um exercício dos sujeitos em níveis variados e ocorre em pontos diferentes da malha social.
Por exemplo, perceber como, nas relações de poder entre homens e \textbf{mulheres}, e entre \textbf{mulheres}, a cultura e a sociedade em nosso tempo histórico constroem diferentes formas de masculinidades e feminilidades.
O que está em jogo, portanto, é rejeitar os determinismos biológicos e dar ênfase ao cultural, nas diferentes relações de poder entre os sexos.
Com isso, usar a perspectiva de gênero é lançar luz de forma interdisciplinar sobre “ a construção social e histórica produzida sobre as características biológicas ” ( LOURO, 1997, p. 22 ).
A presente trilha operacionaliza as seguintes categorias como forma de ferramenta educativa : a ) conhecimento científico : desenvolver a capacidade de compreensão e reflexão com base em conhecimentos legitimados pela comunidade científica ; b ) representações sociais : identificar como são produzidos os modos de ver as \textbf{mulheres}, considerando as interseccionalidades de gênero, as \textbf{populações} LGBTQIAP+, etnia, raça, de classe, geracional, entre outras ; c ) direitos : como campo de emergência e consolidação de direitos da \textbf{mulheres} ; d ) gênero : como categoria e como perspectiva em destaque, na análise da interseccionalidade, ao gerar conhecimentos científicos acerca das \textbf{mulheres} catarinenses ; e ) \textbf{mulheres} : diferenças entre os sujeitos históricos, diferentes em suas lutas ; f ) relações de poder : observar como as práticas sociais legitimam relações assimétricas de poder entre homens e \textbf{mulheres} ; g ) direitos : campo de lutas sociais para garantir igualdade de gênero.
Por fim, a importância desse tema para o Novo Ensino Médio constitui -se na diminuição dessas desigualdades sociais e no reconhecimento das conquistas realizadas pelos movimentos das \textbf{mulheres} e das feministas.
Em parte, os direitos reivindicados foram listados no II Plano Nacional de Políticas para as Mulheres ( PNPM ) ( BRASIL, 2008 ).
Este instrumento definiu metas para implantar políticas públicas de inclusão das questões de gênero, raça e etnia em currículos escolares, com o objetivo de modificar a cultura e a comunicação discriminatórias, por meio de práticas educativas, da produção de conhecimento no âmbito da educação formal.
Nestes termos, o ensino médio enfrenta o desafio maior que é compreender e “ reconhecer a violência de gênero, raça e etnia como violência estrutural e histórica, que expressa a opressão das \textbf{mulheres} e precisa ser tratada como questão de segurança, justiça e saúde pública ” ( BRASIL 2008, p. 28/29 ).
Quadro 17 - Objetos de conhecimento e habilidades da trilha \textbf{mulheres} no \textbf{território} catarinense OBJETOS DE CONHECIMENTO Protagonismos - Representações sociais das \textbf{mulheres} - Masculinidades - Feminilidades - Relações de poder - História dos povos e das \textbf{mulheres} \textbf{indígenas} - Culturas : afrodescendente, afro-brasileira e negra - Patriarcado - Violências de gênero - Mulheres do campo - Diversidade das \textbf{mulheres} no \textbf{território} catarinense - Memórias das \textbf{mulheres} - Direitos humanos e das \textbf{mulheres} - Desigualdades sociais - Direitos das \textbf{mulheres} e ações coletivas - História dos movimentos sociais de \textbf{mulheres} - Movimento feminista - Ondas do feminismo - Machismo e misoginia - Cidadania ativa - Direitos trabalhistas - Formas de violências domésticas, sexuais - Caboclas, \textbf{indígenas}, quilombolas, \textbf{ciganas} - Mulheres e lutas sociais Habilidades - Desnaturalizar, na perspectiva de gênero, as representações sociais das \textbf{mulheres}, no \textbf{território} catarinense. - Identificar \textbf{mulheres} com representatividade comunitária e nos lugares de vivência, inclusive no ambiente escolar. - Historicizar e reconhecer, nas perspectivas sociológicas, históricas, filosóficas e geográficas, os protagonismos das \textbf{mulheres} nos lugares de vivência. - Identificar representações sociais de \textbf{mulheres} com visibilidade pública, tais como atletas, políticas, cientistas, empresárias, professoras, entre outras. - Compreender as representações sociais das masculinidades e das feminilidades em suas expressões marcadas pelas relações de saber/poder, nos lugares de vivência e no ambiente escolar. - Compreender como uma expressão da violência de gênero as representações sociais das \textbf{mulheres} geradas na lógica do patriarcado. - Compreender as representações sociais das \textbf{mulheres} afrodescendentes e \textbf{indígenas} no \textbf{território} catarinense e lugares de vivência. - Identificar e reconhecer as identidades sociais das \textbf{mulheres} caboclas, \textbf{ciganas}, do Contestado, entre outras, na perspectiva de gênero e decolonial no \textbf{território} catarinense. - Reconhecer a história da diversidade étnico-racial e cultural das \textbf{mulheres} \textbf{indígenas} e afro-brasileiras. - Valorizar a cultura e a \textbf{ancestralidade} na diversidade étnico-racial das \textbf{mulheres} \textbf{indígenas} e afro-brasileiras. - Identificar as violências de gênero nas interseccionalidades das \textbf{mulheres} negras, trans, deficientes, entre outras. - Reconhecer os tipos de violência contra as \textbf{mulheres} - física, psicológica, sexual, assédio moral, entre outros - nos lugares de vivência e no ambiente escolar. - Reconhecer as \textbf{mulheres} do campo na construção do espaço rural, na garantia dos direitos à terra, à água, à educação, à saúde, à assistência social, à memória, à história, à segurança, às identidades, entre outros. - Reconhecer os direitos humanos das \textbf{mulheres} com destaque para a igualdade de gênero. - Investigar e reconhecer quais são os direitos das \textbf{mulheres} na contemporaneidade no \textbf{território} catarinense. - Compreender os direitos das \textbf{mulheres} na relação com as desigualdades sociais e as violências de gênero. - Compreender os direitos das \textbf{mulheres} como conquistas das ações coletivas e dos movimentos sociais. - Reconhecer os movimentos ativistas dos coletivos de \textbf{mulheres} no enfrentamento do machismo e da misoginia. - Compreender o que é o machismo na relação igualdade de gênero. - Compreender a misoginia como uma das expressões da violência de gênero. - Historicizar as formas de violências contra as \textbf{mulheres} interseccionadas pelas diversidades étnico-raciais, de gênero e de deficiência, nos lugares de vivência e no ambiente escolar. - Reconhecer as políticas afirmativas e os direitos adquiridos da \textbf{mulheres} \textbf{indígenas} e afro-brasileiras em respeito aos marcos legais no \textbf{território} catarinense. - Reconhecer as memórias das \textbf{mulheres} integrantes dos grupos quilombolas e das \textbf{populações} \textbf{indígenas}, caboclas, \textbf{ciganas}, ribeirinhas, pescadoras, faxinalenses, seringueiras, como patrimônio histórico, no \textbf{território} catarinense. - Identificar os direitos e as legislações das \textbf{mulheres} com deficiência. - Reconhecer a importância dos movimentos feministas no contexto dos movimentos sociais para as garantias dos direitos. - Identificar os direitos das \textbf{mulheres} no \textbf{território} catarinense, considerando as interseccionalidades das identidades e expressões de gênero, raça, etnia, deficiência, entre outros.
Fonte : Elaboração dos autores ( 2020 ).
Quadro 18 – Eixos estruturantes e habilidades da trilha \textbf{mulheres} no \textbf{território} catarinense EIXOS ESTRUTURANTES Investigação científica - Identificar se os direitos das \textbf{mulheres} são garantidos nos lugares de vivência, inclusive no ambiente escolar. - Promover ações que ampliem os direitos das \textbf{mulheres} nos lugares de vivência e no ambiente escolar. - Investigar e visibilizar as desigualdades sociais marcadas por gênero, nos lugares de vivência.
Processos criativos - Propor soluções para problemas em relação aos direitos humanos, combatendo as formas de desigualdade, intolerância, discriminação, injustiça, preconceito e violência contra as \textbf{mulheres} ( machismo, misoginia, racismo, homofobia, xenofobia, entre outros ). - Promover ações que visem à participação plena e ativa das \textbf{mulheres} nas múltiplas esferas da vida pública, inclusive no ambiente escolar.
Mediação e intervenção sociocultural - Propor ações de combate às múltiplas formas de violência contra as \textbf{mulheres} ( física, psicológica, simbólica, sexual, entre outras ), nos lugares de vivência e nas culturas escolares. - Propor ações de reconhecimento dos direitos das \textbf{mulheres}, desnaturalizando as representações das relações sociais nos lugares de vivência. - Propor ações, na escola e nos lugares de vivência, de reconhecimento das culturas e histórias das \textbf{mulheres} afro-brasileiras e \textbf{indígenas} do \textbf{território} catarinense.
Empreendedorismo - Propor ações educativas voltadas aos direitos trabalhistas das \textbf{mulheres}. - Propor ações de reconhecimento da cidadania ativa das \textbf{mulheres}, considerando as interseccionalidades de gênero, as \textbf{populações} LGBTQIAP+, de etnia, raça, classe e geração, entre outras. - Propor ações de reconhecimento de políticas de inclusão social, nas quais se encontram adolescentes, \textbf{mulheres}, idosas, \textbf{mulheres} com deficiência, \textbf{indígenas}, negras, caboclas e quilombolas, \textbf{ciganas}, ribeirinhas, vazanteiras, pescadoras, entre outras \textbf{mulheres} nas suas diversas interseccionalidades. - Oportunizar às \textbf{mulheres} de diferentes gerações a escuta sensível, percebendo a diversidade de trajetórias, de histórias, de memórias e narrativas, nos lugares de vivência.
Fonte : Elaboração dos autores ( 2020 ).

% matched lemmas: ancestralidade, cigano, indígena, mulher, população, racial, sobreviver, território, velho
\end{textsample}
